\begin{afterword}
\summaryhead

Religious fundamentalists, the post says, hold that there can be no morality without a Judge who rewards and punishes: without fear of hell and hope of heaven, nothing would stop people from murdering each other. The author answers with an analogy. If Omega threatened to kill you for using a bathroom in the morning, you would be relieved to see the threat withdrawn. So a believer who is frightened by the thought of God no longer punishing murder must object to murder for reasons of their own. That fear is itself a moral compass, and a believer who loses faith can keep steering by it. Fundamentalists do not ask what would stop people eating pork without hell, which shows that believers do not treat all divine rules alike.

The post then presents two hypothetical philosophers. One argues for selfishness because it benefits society; the other for altruism because it benefits the altruist. Their criteria reveal values opposite to their advice, though which is really the altruist depends on who holds doors for old ladies.

\responsehead

The central argument is short and, for the case it addresses, sound. Someone who is frightened when a threat is withdrawn must have wanted the threatened act prevented anyway, and a threat of hellfire works only because hellfire is already unwelcome. The argument has old relatives. Plato's \textsc{Euthyphro} asks whether the gods love the holy because it is holy; the post asks the psychological version, whether the believer shuns murder because God forbids it. The distinction behind the pork test is drawn in the religious tradition itself: the Talmud separates laws that it would have been logical to write even if they had not been written, bloodshed and theft among them, from statutes with no known reason, the ban on pork among them. That supports the post, and it also shows that the assumption the post ascribes to fundamentalists, that there is ``no moral compass but divine reward and retribution,'' is not a religious view as such.

The argument speaks to the believer who asks the question: if they are frightened, they have a motive of their own, and the post's advice is addressed to that reader. It does not take up William Paley's form of the worry, whether people in general would keep their motivation to do what they know they ought to do without a final sanction. The other reading of the opening claim, that without God nothing would be objectively right or wrong, is not taken up here; the notes on ``Leave a Line of Retreat'' discuss it.

The two philosophers extend the test from fear to argument: the criteria a person appeals to show what counts for them. The post does not rest the verdict on that test. Its last line settles which philosopher is the altruist by what each one does.

In short: a sound argument, with old relatives in Plato and the Talmud, that answers the fundamentalist's question for the believer who asks it.
\end{afterword}
